\scanpage{58}

% 1–2, heading
Construction intrinsèque de $P$ ($k$-lin)

$P\supset K$, et satisfait les conditions

% 3–5, upper part, a)
a) Le foncteur covariant $L\longmapsto\operatorname{Hom}(L,\mathbb E^0)^\vee$ est représentable par un \inserted{objet} $\Psi^0$, Considérons l'hom $\Psi^0\longrightarrow\mathbb E^0$ transposé de

% 6–7, upper part
\[
\begin{aligned}
\operatorname{Hom}(\Psi^0,L)&\longleftarrow\operatorname{Hom}(\mathbb E^0,L)\quad\bigl[\longleftarrow\operatorname{Hom}(\mathbb E^1,L)\\
\rotatebox{90}{$\simeq$}\quad&\hspace{5em}?\\
\operatorname{Hom}(L,\mathbb E^0)^\vee&
\end{aligned}
\]

% 8–10, pairing
déduit de l'accouplement
\[
\begin{aligned}
\operatorname{Hom}(\mathbb E^0,L)\times\operatorname{Hom}(L,\mathbb E^0)&\longrightarrow\operatorname{Hom}(\mathbb E^0,\mathbb E^0)\simeq k\\
(u,\ v)&\longmapsto vu
\end{aligned}
\]

% 11–14, left of argument
\diagram{\begin{tikzcd}[ampersand replacement=\&,column sep=small,row sep=small]
\&\mathbb E^0\arrow[dl,"d"']\arrow[dr,"u"]\&\&\\
\mathbb E^1\arrow[rr,"u'"']\&\&L\arrow[dr,"v"]\&\\
\&\&\&\mathbb E^0
\end{tikzcd}}

% 11–17, centre
Or soit $u'\colon\mathbb E^1\to L$, considérons $u=u'\circ d^0$, et $v\circ u=v\circ(u'\circ d)$, \inserted{$=(v\circ u')\circ d$} comme $vu'=0$, on a donc $v\circ u=0$, donc le composé \deleted{$\operatorname{Hom}$}
\[
\operatorname{Hom}(\Psi^0,L)\longleftarrow\operatorname{Hom}(\mathbb E^0,L)\longleftarrow\operatorname{Hom}(\mathbb E^1,L)
\]
est nul ce qui signifie aussi que
\[
\Psi^0\longrightarrow\mathbb E^0\longrightarrow\mathbb E^1
\]
est nul, donc que $\mathbb E^{\unsure{*|\bullet}}$ est $\Psi^0$-augmenté.

% 16–19, left margin
% slanted
\margin{foncteur $P\to$ \unsure{par|posé|pour} $L\longmapsto\operatorname{Hom}(L,\unsure{\varphi_{\bullet}^{*}|\varphi_{\bullet}^{0}})$}

% 18–19, across middle
Soit $\unsure{\varphi^\bullet|\Phi^\bullet}$ le complexe de cochaînes qu'il forme.

% 19–22, cancelled lines
\deleted{Notons aussi que} \deleted{Soit $P_0$ la sous-catégorie}

\deleted{b) La famille des $\mathbb E^i$ et $\Psi^0$ est génératrice stricte}

\deleted{c) $\varphi^\bullet$ est une résolution de}

% 23–26, above b)
$k$-lin. \deleted{engendrée} formée des sommes finies d'objets de la forme $\mathbb E^i,\varphi_{\bullet}^{0}$, i.e. les expressions de la forme
\[
L+(\Psi^0)^r\qquad r\in\mathbb N,\ L^{\unsure{\,|'}}\in\operatorname{Ob} K
\]

% 27, at left and in margin
\inserted{On a}
% slanted
\margin{plus bas \unsure{voir|si} \unsure{ce qu'il} \unsure{faudrait}}

% 27–30, rewritten b)
b) Les suites $0\to L'\to L$ \deleted{\unsure{telles que}} \inserted{$0\to L'\to L\to L''$,} \unsure{$0\to L'\to L\to L''\to0,\ L'\to L\to L''\to0$|$\to L''\to0,\ L'\to L\to L''\to0$}, telles que les $P\in P_0$ sont projectifs \unsure{pour celles|pour celles-ci|fournit}, i.e. telles que les $0\to L'\to L\to L''\to0$ soient exactes pour \unsure{les dits|les dits $P$}, sont exactes.

% 31–33, foot
\deleted{\unsure{c)}} Soit $\Sigma$ la classe de ces suites.

\deleted{\unsure{b)}}c) Le complexe augmenté $(\mathbb E^\bullet,\Psi^0,\varepsilon)$ est une $\Sigma$-résolution, i.e. $\varphi^\bullet$ est $\Sigma$-acyclique

% archive, bottom left
% archive sheet number: 57
